\begin{spacing}{1.5}
\vspace{1cm}

\textit{The implementation chapter explains how your proposed approach (from Chapter 3) was actually built: the tools and technologies used, the system's architecture, and the details of how each part was implemented. Replace all italicized instruction text in this chapter with your own writing, and delete this paragraph.}

\section{Overview}

\textit{Briefly describe what was implemented and outline what each section of this chapter covers.}

\section{Technologies Used}

\textit{List the programming languages, frameworks, libraries, hardware platforms, and tools used for implementation, including version numbers where relevant. A short table works well here.}

\section{Dataset Preparation}

\textit{Describe how the dataset (from Chapter 3) was prepared for use in the implementation: file formats, directory structure, train/validation/test splits, augmentation pipelines, etc. Remove this section if not applicable to your project.}

\subsection{Class Definitions}

\textit{If your project involves classification or categorization, define and describe each class/category here. Add code listings as needed, for example:}

% \begin{lstlisting}[language=Python]
% def example_function():
%     pass
% \end{lstlisting}

\section{System Architecture}

\textit{Describe the overall architecture of the system you implemented. Include diagrams, flowcharts, block diagrams, or any visual aids that help the reader understand the system's structure and how components interact.}

\section{Implementation Details}

\textit{Walk through how the system/model was actually implemented, module by module or stage by stage. Include relevant code snippets, configuration details, hyperparameters, and explanations of key design decisions. Discuss any significant challenges encountered and how they were resolved.}

\end{spacing}
